\documentclass[lettersize,journal]{IEEEtran}

\usepackage{amsmath,amsfonts,amssymb}
\usepackage{algorithmic}
\usepackage{algorithm}
\usepackage{array}
\usepackage[caption=false,font=normalsize,labelfont=sf,textfont=sf]{subfig}
\usepackage{textcomp}
\usepackage{stfloats}
\usepackage{url}
\usepackage{verbatim}
\usepackage{graphicx}
\usepackage{lineno,hyperref}
\usepackage{cite}
\usepackage{orcidlink}

\hyphenation{op-tical net-works semi-conduc-tor IEEE-Xplore}

\begin{document}

\title{HDRNet: Hierarchical Damage-aware Relation Refinement Network for Post-Disaster UAV Image Segmentation}

\author{Mingjie~Zhang~\orcidlink{0009-0001-8666-6879},
        Kai~Chen~\orcidlink{0000-0002-1145-0117},
        Yongxu~Zhou~\orcidlink{0009-0007-3464-6714},
        and~Weifeng~Shan~\orcidlink{0000-0002-9658-9483}
\thanks{This work was funded by the Hebei Natural Science Foundation (D2026512004) and the Scientific Research and Development Plan Project of Langfang Science and Technology Bureau (2025011043).}
\thanks{Mingjie Zhang, Kai Chen, and Yongxu Zhou are with the School of Information Science and Technology, Beijing University of Technology, Beijing 100811, China; Beijing Laboratory of Smart Environmental Protection, China; and Beijing Artificial Intelligence Institute, China (e-mail: zhangmingjie@emails.bjut.edu.cn; myck919@emails.bjut.edu.cn; bhhtzyx@nciae.edu.cn).}
\thanks{Weifeng Shan is with the School of Emergency Management, University of Emergency Management, Sanhe 065201, China, and Hebei Key Laboratory of Resource and Environmental Disaster Mechanism and Risk Monitoring, Sanhe 065201, China (e-mail: william.shan@gmail.com).}
\thanks{Corresponding author: Weifeng Shan.}}

\markboth{IEEE Transactions on Geoscience and Remote Sensing,~Vol.~XX, No.~XX, Month~Year}%
{Zhang \MakeLowercase{\textit{et al.}}: HDRNet for Post-Disaster UAV Image Segmentation}

\maketitle

\begin{abstract}
Semantic segmentation of post-disaster unmanned aerial vehicle (UAV) imagery is a key enabling technology for rapid emergency assessment, underpinning tasks such as flood monitoring, road accessibility analysis, building damage evaluation, and waterfront infrastructure inspection. Relative to ordinary scenes, post-disaster UAV imagery raises three intertwined challenges: strong contextual dependencies among disaster-related objects, fine-grained and hierarchically organized damage categories that are easily confused, and fragmented or incomplete object boundaries. To tackle these issues within a unified framework, we propose HDRNet, a hierarchical damage-aware relation refinement network built upon a SegFormer encoder. First, a Disaster Relation-aware Context Aggregation (DRCA) module adaptively aggregates multi-scale context conditioned on coarse disaster semantics and incorporates category co-occurrence priors to model the relationships among flooded areas, buildings, roads, vegetation, vehicles, and damaged infrastructure. Second, a Hierarchical Damage-aware Decoding (HDAD) head decouples segmentation into object-level recognition and disaster-state estimation, offering a unified formulation for flooded/non-flooded, damaged/intact, and multi-level destruction labels. Third, a Damage-sensitive Structural Boundary Refinement (DSBR) module combines automatically generated boundary supervision with learnable structural responses to refine fragmented contours in submerged roads, damaged buildings, blocked roads, and waterfront structures. Experiments on three public disaster-scene datasets, FloodNet, RescueNet, and FWISD, show that HDRNet consistently outperforms representative CNN-, Transformer-, and state-space-based methods, with especially clear gains on damage-related categories.
\end{abstract}

\begin{IEEEkeywords}
Semantic segmentation, post-disaster unmanned aerial vehicle (UAV) imagery, damage assessment, contextual relation modeling, hierarchical damage decoding, boundary refinement.
\end{IEEEkeywords}

\section{Introduction}
\IEEEPARstart{N}{atural} disasters such as hurricanes and floods pose growing threats to human lives and infrastructure, making rapid and accurate post-disaster scene understanding indispensable for emergency response and damage assessment \cite{liuStairFusionNetwork2024, liuClusterformerPineTree2024, sarkarSAMVQASupervisedAttentionBased2023}.
Unmanned aerial vehicles (UAVs) have emerged as a vital platform for post-disaster data collection, since they can rapidly acquire high-resolution imagery over large and otherwise inaccessible areas\cite{luLightweightCNNTransformer2024, liuLightWeightSemanticSegmentation2021,salehLiSTNetEnhancedFlood2024}.
Semantic segmentation of such UAV imagery yields a pixel-level interpretation of the scene, directly supporting critical tasks including flood monitoring, road accessibility analysis, building damage evaluation, and waterfront infrastructure inspection\cite{mouRelationMattersRelational2020, niuHybridMultipleAttention2022,piDetectionSemanticSegmentation2021}.

Compared with ordinary natural scenes or conventional remote-sensing imagery, post-disaster UAV images are considerably more challenging, as the task is inherently fine-grained and damage-oriented rather than a mere categorization of object types. Through an analysis of three representative public datasets, FloodNet\cite{rahnemoonfarFloodnetHighResolution2021a}, RescueNet\cite{rahnemoonfarRescueNetHighResolution2023a}, and FWISD\cite{xueFWISDFloodWaterfront2026a}, we identify three intertwined challenges that current segmentation methods do not explicitly address.

First, disaster-related objects exhibit strong contextual dependencies rather than appearing in isolation. Flooded roads typically lie adjacent to large bodies of water, flooded buildings appear near inundated regions, blocked roads co-occur with damaged buildings and debris, and damaged waterfront infrastructure is spatially coupled with shorelines and water. Visually similar categories, such as floodwater versus pools or shadowed roads versus buildings, can be disambiguated only from their surrounding context. Conventional context modules such as ASPP, PPM, or generic self-attention enlarge the receptive field but fail to explicitly model the semantic co-occurrence and spatial dependency among disaster categories\cite{chenRethinkingAtrousConvolution2017, liuStairFusionNetwork2024, zhaoPyramidSceneParsing2017}.

Second, disaster categories are hierarchically structured rather than flat. Object types and their damage states are naturally coupled: a building may be intact or destroyed at several levels (no damage, minor, major, or total destruction), and a road may be clear or blocked. Because adjacent damage levels differ only subtly in appearance, directly applying a flat softmax over all leaf categories tends to confuse neighboring damage grades and disregards the underlying object--state hierarchy\cite{fuDeepOrdinalRegression2018a, liDeepHierarchicalSemantic2022}.

Third, disaster-related objects often exhibit fragmented and incomplete boundaries. Submerged roads, damaged buildings, blocked roads, and thin elongated waterfront structures frequently show blurred or broken contours owing to occlusion, inundation, and shadows. Standard decoders that upsample low-resolution semantic features tend to oversmooth these boundaries, which is particularly detrimental for small or slender disaster targets whose accurate delineation is crucial for assessment\cite{kervadecBoundaryLossHighly2021, takikawaGatedSCNNGatedShape2019}.

Existing UAV and remote-sensing segmentation methods, including recent Transformer- and state-space-based models, primarily pursue stronger generic feature representations and seldom address these three disaster-specific issues within a unified framework\cite{wangUAVSegDualEncoderCrossScale2025,maMultilevelMultimodalFusion2024,huangMambaUAVSegNetMultiscaleAdaptive2024}. 

Consequently, they underperform precisely on the most decision-critical categories, namely flooded roads, multi-level building damage, blocked roads, and damaged waterfront infrastructure.

To address these challenges, we recast the problem as fine-grained damage-aware semantic segmentation of post-disaster UAV imagery and propose HDRNet, a Hierarchical Damage-aware Relation Refinement Network built upon a SegFormer encoder. HDRNet integrates three task-specific modules, each targeting one of the challenges above. A Disaster Relation-aware Context Aggregation (DRCA) module adaptively aggregates multi-scale context conditioned on coarse disaster semantics and injects category co-occurrence priors to model the relationships among flooded areas, buildings, roads, vegetation, vehicles, and infrastructure. A Hierarchical Damage-aware Decoding (HDAD) head decouples segmentation into object-level recognition and disaster-state estimation, yielding a unified formulation for flooded/non-flooded, damaged/intact, and multi-level destruction labels across datasets. A Damage-sensitive Structural Boundary Refinement (DSBR) module leverages automatically generated boundary supervision and learnable structural responses to sharpen fragmented disaster boundaries. The overall design follows a coherent pipeline of global relation modeling, hierarchical damage decoding, and local structural refinement.

The main contributions of this work are summarized as follows:
\begin{itemize}
\item We propose HDRNet, a hierarchical damage-aware relation refinement network for UAV-based post-disaster semantic segmentation that jointly addresses contextual ambiguity, fine-grained damage confusion, and boundary fragmentation within a single framework.
\item We design a Disaster Relation-aware Context Aggregation (DRCA) module that adaptively aggregates multi-scale context conditioned on coarse disaster semantics and explicitly models category co-occurrence relationships among disaster-related objects.
\item We introduce a Hierarchical Damage-aware Decoding (HDAD) head that decomposes disaster-scene segmentation into object-level recognition and disaster-state estimation, providing a unified formulation for flooded/non-flooded, damaged/intact, and multi-level destruction categories across FloodNet, RescueNet, and FWISD.
\item We develop a Damage-sensitive Structural Boundary Refinement (DSBR) module that employs automatically generated boundary supervision and learnable structural responses to refine the boundaries of submerged roads, damaged buildings, blocked roads, and waterfront infrastructure.
\end{itemize}

Extensive experiments on FloodNet, RescueNet, and FWISD show that HDRNet consistently outperforms representative CNN-, Transformer-, and state-space-based segmentation methods, with the largest gains on the most challenging damage-related categories. The remainder of this paper is organized as follows. Section~II reviews related work; Section~III details the proposed HDRNet; Section~IV reports experiments and analyses; and Section~V concludes the paper.

\section{Related Work}

\subsection{Semantic Segmentation}
Deep learning has driven rapid progress in semantic segmentation. Early fully convolutional networks established the paradigm of dense pixel-wise prediction, after which encoder--decoder architectures, dilated/atrous convolutions, and multi-scale context modules such as PPM and ASPP became standard tools for capturing global context\cite{longFullyConvolutionalNetworks2015,ronnebergerUNetConvolutionalNetworks2015,chenEncoderDecoderAtrousSeparable2018,zhaoPyramidSceneParsing2017}.
Transformer-based models have since reshaped the field: vision Transformers and hierarchical variants such as SegFormer deliver strong long-range modeling with efficient multi-scale feature extraction, whereas query-based formulations such as K-Net and Mask2Former unify semantic and instance segmentation under a single framework\cite{dosovitskiyImageWorth16x162020,xie2021segformer,zhang2021knet,NEURIPS2021_950a4152}. 
More recently, state-space and large-kernel models have emerged as efficient alternatives to attention for high-resolution imagery, and foundation-model-based formulations have further broadened the design space of remote-sensing segmentation\cite{lou2025overlock, SegMAN, xiao2025spatialmamba, yeExploitingFoundationModels2026, biViRefSAMVisualReference2026}. 
These methods primarily pursue stronger and more efficient generic representations. However, they are tailored to general scenes and do not explicitly exploit the task-specific structure of post-disaster imagery, in which contextual relations, damage hierarchy, and fragmented boundaries play a central role. In this work, we adopt SegFormer as the backbone and build disaster-specific modules upon it.

\subsection{UAV and Disaster-Scene Segmentation}
Owing to their flexibility and high spatial resolution, UAVs have become a key platform for post-disaster data acquisition\cite{liuClusterformerPineTree2024,shaheenSwinSegFormerAdvancingAerial2025}.
To support data-driven research, several high-resolution UAV benchmarks have been released, including FloodNet for post-flood scene understanding, RescueNet for natural-disaster damage assessment, and FWISD for flood and waterfront infrastructure damage\cite{rahnemoonfarFloodnetHighResolution2021, rahnemoonfarRescueNetHighResolution2023, xueFWISDFloodWaterfront2026}. 
Building on such data, numerous UAV- and remote-sensing-oriented segmentation methods have been proposed, ranging from lightweight real-time networks and dual-encoder cross-scale designs to multi-level multimodal fusion Transformers and state-space backbones tailored to aerial imagery\cite{wangUAVSegDualEncoderCrossScale2025,maMultilevelMultimodalFusion2024,huangMambaUAVSegNetMultiscaleAdaptive2024,shaheenSwinSegFormerAdvancingAerial2025,chenRSMambaRemoteSensing2024,felirenProgressiveCrossAttention2025}.
Several recent studies further target disaster scenarios specifically, including detection-guided and damage-aware segmentation of UAV footage\cite{piDetectionSemanticSegmentation2021,zetoutCSDNetContextAwareSegmentation2025,zhuDASegFormerDamageAwareSemantic2026}.
Nevertheless, most of these methods treat disaster categories as flat classes and concentrate on generic feature enhancement, leaving the coupling between object identity and damage state, as well as the disambiguation of contextually similar categories, largely unaddressed. In contrast, our work explicitly formulates disaster segmentation as a fine-grained, damage-aware task and evaluates it across three datasets.

\subsection{Context Modeling, Hierarchical Labeling, and Boundary Refinement}
The three challenges addressed in this paper correspond to three lines of research. First, context modeling has been studied extensively through multi-scale pooling, atrous convolution, and self-attention; these techniques enlarge the receptive field but typically rely on fixed aggregation and do not encode category co-occurrence priors. In contrast, our DRCA module aggregates context adaptively according to coarse disaster semantics and injects explicit relation priors\cite{maRS3MambaVisualState2024, chenAfaMambaAdaptiveFeature2025, yuanObjectContextualRepresentationsSemantic2020}. 
Second, hierarchical and fine-grained labeling has been explored to model label structure and ordinal relations among classes, yet it is rarely tailored to the object--damage hierarchy of disaster scenes; our HDAD head decouples object recognition from disaster-state estimation to reduce confusion among adjacent damage levels\cite{fuDeepOrdinalRegression2018a, liDeepHierarchicalSemantic2022a, weberFlatteningParentBias2024}. 
Third, boundary-aware segmentation improves contour quality through boundary supervision or contour-oriented losses, yet generic boundary modules are insensitive to the fragmented structures typical of damaged targets; our DSBR module leverages automatically generated boundary supervision and learnable structural responses to refine disaster-related boundaries\cite{kervadecBoundaryLossHighly2021a, takikawaGatedSCNNGatedShape2019a, zhangBoundaryAwareSpatialFrequency2024, yuProbabilityGuidedEdgeEnhancement2025, sunHandlingNoisyAnnotation2025}.
By integrating these three task-specific designs into a single SegFormer-based framework, HDRNet offers a unified solution to challenges that existing approaches address only partially.

\section{Methodology}

\subsection{Overall Architecture}
The overall architecture of the proposed HDRNet is illustrated in Fig.~\ref{fig:overview}. HDRNet employs SegFormer-B2 (MiT-B2) as the encoder, which extracts four multi-scale features from an input image $I\in\mathbb{R}^{H\times W\times 3}$:
\begin{equation}
F_1, F_2, F_3, F_4 = E(I),
\end{equation}
where $F_1$ is a shallow high-resolution feature carrying edges and local structures, $F_2$ and $F_3$ encode mid- and high-level structural and object semantics, and $F_4$ is the deepest feature capturing global context. The Disaster Relation-aware Context Aggregation (DRCA) module operates on the high-level feature to model disaster-specific contextual relations:
\begin{equation}
\hat{F}_4 = \mathrm{DRCA}(F_4).
\end{equation}
A lightweight MLP decoder then fuses the multi-scale features into a decoded representation:
\begin{equation}
F_d = D(F_1, F_2, F_3, \hat{F}_4).
\end{equation}
The Damage-sensitive Structural Boundary Refinement (DSBR) module subsequently refines the decoded feature using low-level structural cues, and the Hierarchical Damage-aware Decoding (HDAD) head produces the final prediction together with auxiliary object and damage-state predictions:
\begin{equation}
F_r = \mathrm{DSBR}(F_d, F_1, F_2),
\end{equation}
\begin{equation}
P, P_o, P_d = \mathrm{HDAD}(F_r),
\end{equation}
where $P$ is the final semantic segmentation, $P_o$ is the object-category prediction, and $P_d$ is the disaster-state prediction. Together, the three modules form a coherent pipeline of global relation modeling, hierarchical damage decoding, and local structural refinement, each addressing one of the three challenges identified in Section~I.

\begin{figure*}[!t]
\centering
% TODO: replace with the HDRNet overall-architecture figure (file: fig/method_overview.pdf)
\IfFileExists{fig/method_overview.pdf}
  {\includegraphics[width=0.95\textwidth]{fig/method_overview.pdf}}
  {\fbox{\parbox[c][0.28\textheight][c]{0.93\textwidth}{\centering Placeholder: overall architecture of HDRNet (DRCA $\rightarrow$ decoder $\rightarrow$ DSBR $\rightarrow$ HDAD). Replace with \texttt{fig/method\_overview.pdf}.}}}
\caption{Overall architecture of the proposed HDRNet. The SegFormer-B2 encoder extracts multi-scale features $F_1$--$F_4$; DRCA enhances the high-level feature with disaster relation-aware context; an MLP decoder fuses multi-scale features; DSBR refines boundaries using low-level structural cues; and the HDAD head jointly predicts the object category, the disaster state, and the final segmentation map.}
\label{fig:overview}
\end{figure*}

\subsection{Disaster Relation-aware Context Aggregation}
Post-disaster scenes exhibit strong contextual dependencies among categories: flooded roads co-occur with water, submerged buildings are surrounded by inundated regions, and blocked roads appear near collapsed structures. Conventional context modules such as ASPP and PPM aggregate multi-scale information with fixed receptive fields, while generic global self-attention fails to explicitly encode the semantic co-occurrence among disaster categories. The proposed Disaster Relation-aware Context Aggregation (DRCA) module addresses this limitation through three components: a multi-scale context branch, a coarse disaster-semantic gating mechanism, and a category relation prior.

\textit{Multi-scale context branch.} Given the high-level feature $F_4$, DRCA constructs $N$ parallel branches to capture disaster objects at different scales. A pointwise branch $\mathrm{Conv}_{1\times1}(F)$ preserves local channel information; three depthwise branches $\mathrm{DWConv}_{3\times3}$ with dilation rates $1$, $6$, and $12$ capture small structures (vehicles), mid-scale structures (roads, buildings), and large regions (water bodies, flooded areas), respectively; and a global branch $\mathrm{Broadcast}(\mathrm{MLP}(\mathrm{GAP}(F)))$ injects image-level context. Depthwise separable convolutions are adopted to keep the computational overhead low on high-resolution UAV imagery. The branch outputs are denoted $\{F_1^s,\dots,F_N^s\}$.

\textit{Coarse disaster-semantic gating.} Rather than fusing the branches with fixed weights, DRCA generates the fusion weights adaptively from the coarse disaster semantics of each image. An image-level descriptor $z=\mathrm{GAP}(F)$ first predicts a coarse disaster-semantic distribution
\begin{equation}
p_c = \mathrm{Softmax}(H_c(z)),
\end{equation}
over coarse categories (water/flood, building/infrastructure, road, vegetation, vehicle/others). The branch weights are then obtained as
\begin{equation}
\alpha = \mathrm{Softmax}(W_g\,p_c),\quad \alpha=[\alpha_1,\dots,\alpha_N],
\end{equation}
and the gated multi-scale context is
\begin{equation}
F_{ctx} = \sum_{i=1}^{N}\alpha_i\,F_i^s.
\end{equation}
This enables the model to emphasize large-scale and global branches when water or flooded regions dominate, and local and mid-scale branches when roads, vehicles, or building details prevail.

\textit{Category relation prior.} To further strengthen relational reasoning among disaster categories, DRCA injects a co-occurrence prior estimated from the training set. The conditional co-occurrence of category $c_j$ given category $c_i$ is
\begin{equation}
R_{ij} = \frac{N(c_i,c_j)}{N(c_i)},
\end{equation}
where $N(\cdot)$ counts (co-)occurrences over images. The prior is incorporated as an attention bias during relation modeling,
\begin{equation}
A_{ij} = \mathrm{Softmax}\!\left(\frac{Q_iK_j^{\top}}{\sqrt{d}} + \beta R_{ij}\right),
\end{equation}
where $Q_i$ and $K_j$ are query and key features, $d$ is the feature dimension, and $\beta$ controls the prior strength. The relation-enhanced context is $\hat{F}_{ctx}=F_{ctx}+\mathrm{RelationAttn}(F_{ctx},R)$, and a final $1\times1$ projection yields $\hat{F}_4=\mathrm{Conv}_{1\times1}(\hat{F}_{ctx})$.

\begin{figure}[!t]
\centering
% TODO: replace with the DRCA module figure (file: fig/drca.pdf)
\IfFileExists{fig/drca.pdf}
  {\includegraphics[width=0.95\linewidth]{fig/drca.pdf}}
  {\fbox{\parbox[c][0.22\textheight][c]{0.9\linewidth}{\centering Placeholder: DRCA module (multi-scale branches $\rightarrow$ coarse-semantic gating $\rightarrow$ relation prior). Replace with \texttt{fig/drca.pdf}.}}}
\caption{Structure of the Disaster Relation-aware Context Aggregation (DRCA) module. Multi-scale context branches are adaptively fused by coarse disaster-semantic gating, and a category co-occurrence prior is injected as an attention bias to model relations among disaster categories.}
\label{fig:drca}
\end{figure}

\subsection{Hierarchical Damage-aware Decoding}
Disaster-scene categories are inherently hierarchical: a label such as \emph{flooded building} couples an object identity (building) with a disaster state (flooded), and RescueNet further distinguishes ordered damage levels (no/minor/major/total destruction). Treating such composite labels as flat categories ignores their shared structure and tends to confuse adjacent damage levels. The Hierarchical Damage-aware Decoding (HDAD) head therefore decomposes the prediction into an object branch, a disaster-state branch, and a fused final branch, coupling them through a hierarchical consistency constraint.

\textit{Object and disaster-state branches.} Given the decoded feature $F_d$, the object branch predicts base object categories (e.g., building, road, water, vegetation, vehicle), and the disaster-state branch predicts the damage status:
\begin{equation}
P_o = \mathrm{Softmax}(H_o(F_d)),\qquad P_d = \mathrm{Softmax}(H_d(F_d)).
\end{equation}
The corresponding supervisions $Y_o$ and $Y_d$ are derived from the original labels by mapping composite classes to their object and state factors (e.g., flooded/non-flooded building $\rightarrow$ building; minor/major/total destruction $\rightarrow$ building). For the ordered building-damage levels in RescueNet, the $K$-class problem is cast as $K-1$ binary decisions via ordinal regression,
\begin{equation}
P(y\geq k\,|\,f) = \sigma(w_k^{\top}f + b_k),\quad k=1,\dots,K-1,
\end{equation}
with $P(y=k)=P(y\geq k)-P(y\geq k+1)$ and boundary conditions $P(y\geq 0)=1$ and $P(y\geq K)=0$. The ordinal loss is $\mathcal{L}_{ord}=\sum_{k=1}^{K-1}\mathrm{BCE}(P(y\geq k),Y_k^{ord})$.

\textit{Fused final branch.} The final semantic prediction fuses the decoded feature with the object and damage features:
\begin{equation}
F_f = \mathrm{Fuse}([F_d, F_o, F_{dmg}]),\qquad P = \mathrm{Softmax}(H_f(F_f)),
\end{equation}
where $F_o=H_o^{feat}(F_d)$ and $F_{dmg}=H_d^{feat}(F_d)$.

\textit{Hierarchical consistency.} To keep the three predictions mutually consistent, mapping matrices $M_o$ and $M_d$ project the final prediction onto implicit object and state distributions, $\tilde{P}_o=M_oP$ and $\tilde{P}_d=M_dP$. The consistency loss penalizes their divergence from the dedicated branches,
\begin{equation}
\mathcal{L}_{hier} = D_{KL}(P_o\,\|\,\tilde{P}_o) + D_{KL}(P_d\,\|\,\tilde{P}_d).
\end{equation}
This unified formulation accommodates flooded/non-flooded (FloodNet), damaged/intact (FWISD), and multi-level destruction (RescueNet) labels within a single head, while reducing semantic conflicts between the final prediction and its hierarchical factors.

\subsection{Damage-sensitive Structural Boundary Refinement}
Many disaster-relevant categories suffer from blurred and fragmented boundaries: the flooded-road/water interface is ambiguous, submerged buildings are absorbed by water and shadow, damaged buildings exhibit irregular contours, and waterfront infrastructure appears as thin structures easily confused with the background. Standard decoders prioritize semantic recovery and tend to over-smooth these boundaries during upsampling. The Damage-sensitive Structural Boundary Refinement (DSBR) module introduces explicit boundary supervision and a structural-response branch to sharpen them.

\textit{Boundary supervision without extra labels.} Boundary targets are derived automatically from the semantic mask through a morphological gradient,
\begin{equation}
Y_b = \mathrm{Dilate}(Y) - \mathrm{Erode}(Y),
\end{equation}
implemented with max-pooling for dilation and negative max-pooling for erosion. To emphasize the task, boundaries can be generated only for disaster-relevant classes (e.g., flooded building/road and water for FloodNet; damaged buildings and blocked roads for RescueNet; damaged and waterfront structures for FWISD). An auxiliary boundary prediction is then produced from a low-level feature, $P_b=\sigma(H_b(F_l))$.

\textit{Structural-response branch.} A structural branch $F_{str}=H_{str}(F_l)$ enhances directionally structured targets such as roads, buildings, and infrastructure. Its convolutions are initialized with horizontal, vertical, diagonal, and anti-diagonal kernels but remain learnable, so that they can adapt to deformed, occluded, and fragmented post-disaster structures.

\textit{Boundary attention injection.} The boundary and structural cues are fused into a boundary attention map that re-weights the decoded feature,
\begin{equation}
M_b = \sigma(\mathrm{Conv}([P_b, F_{str}])),\qquad \hat{F}_d = F_d + F_d \odot M_b,
\end{equation}
where $\odot$ denotes element-wise multiplication, and the refined feature is $F_r=\mathrm{Conv}_{1\times1}(\hat{F}_d)$. The boundary prediction is supervised by a combination of BCE and Dice losses, $\mathcal{L}_{boundary}=\mathrm{BCE}(P_b,Y_b)+\mathrm{Dice}(P_b,Y_b)$.

\begin{figure}[!t]
\centering
% TODO: replace with the HDAD + DSBR figure (file: fig/hdad_dsbr.pdf)
\IfFileExists{fig/hdad_dsbr.pdf}
  {\includegraphics[width=0.95\linewidth]{fig/hdad_dsbr.pdf}}
  {\fbox{\parbox[c][0.24\textheight][c]{0.9\linewidth}{\centering Placeholder: HDAD decoding head (object / disaster-state / fused branches with hierarchical consistency) and DSBR boundary refinement. Replace with \texttt{fig/hdad\_dsbr.pdf}.}}}
\caption{The Hierarchical Damage-aware Decoding (HDAD) head and the Damage-sensitive Structural Boundary Refinement (DSBR) module. HDAD jointly predicts object categories, disaster states, and the fused final segmentation under a hierarchical consistency constraint, while DSBR injects boundary attention derived from automatically generated boundary supervision and learnable structural responses.}
\label{fig:hdad-dsbr}
\end{figure}

\subsection{Training Objective}
HDRNet is trained end-to-end with a multi-task objective that combines the main segmentation loss with the coarse-semantic, object, disaster-state, hierarchical-consistency, and boundary losses:
\begin{equation}
\begin{split}
\mathcal{L}_{total} = {}& \mathcal{L}_{seg} + \lambda_1\mathcal{L}_{coarse} + \lambda_2\mathcal{L}_{obj} \\
&+ \lambda_3\mathcal{L}_{damage} + \lambda_4\mathcal{L}_{hier} + \lambda_5\mathcal{L}_{boundary}.
\end{split}
\end{equation}
The main segmentation loss combines cross-entropy and Dice terms, $\mathcal{L}_{seg}=\mathrm{CE}(P,Y)+\mathrm{Dice}(P,Y)$. The coarse-semantic loss $\mathcal{L}_{coarse}=\mathrm{CE}(P_c,Y_c)$ supervises the DRCA gating distribution; the object and disaster-state losses are $\mathcal{L}_{obj}=\mathrm{CE}(P_o,Y_o)$ and $\mathcal{L}_{damage}=\mathrm{CE}(P_d,Y_d)$ (the latter replaced by the ordinal loss $\mathcal{L}_{ord}$ for RescueNet damage levels); and $\mathcal{L}_{hier}$ and $\mathcal{L}_{boundary}$ are defined above. Unless otherwise stated, the loss weights are set to $\lambda_1=0.1$, $\lambda_2=0.4$, $\lambda_3=0.6$, $\lambda_4=0.3$, and $\lambda_5=0.4$.

\section{Experiments}
\subsection{Data Description}
We evaluate our method on three public benchmarks for disaster-scene semantic segmentation: FloodNet\cite{rahnemoonfarFloodnetHighResolution2021a}, RescueNet\cite{rahnemoonfarRescueNetHighResolution2023a}, and FWISD\cite{xueFWISDFloodWaterfront2026a}. Each dataset comprises post-disaster UAV imagery with pixel-level annotations and poses challenging conditions for segmenting flooded areas, damaged buildings, roads, water bodies, vegetation, vehicles, and other disaster-related objects. Representative images and their pixel-level annotations are shown in Fig.~\ref{fig:dataset-examples}, together with the class color legend used throughout this paper. The datasets are detailed below.

\begin{figure*}[!t]
\centering
\includegraphics[width=0.72\textwidth]{fig/dataset_examples.pdf}
\caption{Representative examples from the three disaster-scene benchmarks. For each dataset, two UAV images are shown alongside their ground-truth semantic annotations, with the per-class color legend on the right. The same color mapping is used for all qualitative results in this paper.}
\label{fig:dataset-examples}
\end{figure*}

1) \textit{FloodNet}: FloodNet is a high-resolution UAV dataset acquired in the aftermath of Hurricane Harvey for post-flood scene understanding. It supplies pixel-wise semantic labels for disaster-related categories, including flooded and non-flooded buildings, flooded and non-flooded roads, water, trees, vehicles, pools, grass, and background. Because the original images are of large spatial size, we crop them into $1024 \times 1024$ patches for training and evaluation. After preprocessing, the FloodNet subset used in this paper comprises 17,992 training images and 449 testing images.

2) \textit{RescueNet}: RescueNet is a high-resolution UAV segmentation dataset tailored to natural-disaster damage assessment. Collected over regions struck by Hurricane Michael, it offers dense semantic labels for scene elements such as water, buildings, vehicles, roads, trees, and pools. Notably, RescueNet defines fine-grained damage-related categories, including multiple building damage levels and road conditions, which makes it well suited to evaluating segmentation models in complex post-disaster environments. As with FloodNet, the original high-resolution images are cropped into $1024 \times 1024$ patches. The processed RescueNet subset contains 43,143 training images and 451 testing images.

3) \textit{FWISD}: FWISD is built from high-resolution UAV imagery for flood and waterfront infrastructure damage assessment. It targets the semantic segmentation of hurricane-induced damage to critical waterfront structures and provides labels for 11 classes that distinguish intact from damaged infrastructure-related objects. Since the original FWISD samples are already supplied at the target scale, we retain their native resolution without further cropping. The FWISD subset used in our experiments comprises 2,771 training images and 981 testing images.

\subsection{Implementation Details}
All experiments are implemented within the MMSegmentation framework. Unless otherwise stated, the proposed HDRNet employs SegFormer-B2 (MiT-B2) as the backbone, initialized with ImageNet pre-trained weights, so that the observed gains can be attributed to the proposed modules rather than to a heavier encoder. For fair comparison, all competing methods are trained under the identical data pipeline and schedule. During training, the input images are randomly resized within a ratio range of $(0.5, 2.0)$ and cropped to $512 \times 512$, with the maximum single-category ratio per crop capped at $0.75$ to mitigate class imbalance. We further adopt random horizontal and vertical flipping, random rotation, and photometric distortion for data augmentation. The models are optimized with AdamW using an initial learning rate of $6\times10^{-5}$, a weight decay of $0.01$, and $(\beta_1,\beta_2)=(0.9,0.999)$, under a polynomial learning-rate decay with power $1.0$. Each model is trained for $80{,}000$ iterations on two GPUs with a batch size of $8$ per GPU (total batch size $16$) and evaluated every $10{,}000$ iterations. The loss weights are set to $\lambda_1=0.1$, $\lambda_2=0.4$, $\lambda_3=0.6$, $\lambda_4=0.3$, and $\lambda_5=0.4$ for the coarse-semantic, object, damage, hierarchy-consistency, and boundary terms, respectively. At inference, we employ sliding-window prediction with a crop size of $512 \times 512$ and a stride of $384 \times 384$, yielding a $128$-pixel overlap between adjacent windows to suppress border artifacts.

\subsection{Evaluation Metrics}
Following common practice in semantic segmentation, we report the per-class Intersection-over-Union (IoU) and F1-score, along with their dataset-level averages, namely the mean IoU (mIoU) and mean F1-score (mF1). For class $c$, let $\mathrm{TP}_c$, $\mathrm{FP}_c$, and $\mathrm{FN}_c$ denote the numbers of true-positive, false-positive, and false-negative pixels. The IoU and F1-score are defined as
\begin{equation}
\begin{aligned}
\mathrm{IoU}_c&=\frac{\mathrm{TP}_c}{\mathrm{TP}_c+\mathrm{FP}_c+\mathrm{FN}_c},\\
\mathrm{F1}_c&=\frac{2\,\mathrm{TP}_c}{2\,\mathrm{TP}_c+\mathrm{FP}_c+\mathrm{FN}_c}.
\end{aligned}
\end{equation}
The mIoU and mF1 are obtained by averaging $\mathrm{IoU}_c$ and $\mathrm{F1}_c$ over all classes, and mIoU serves as the primary metric for ranking the methods. To better characterize each model's behavior on disaster-related objects, we place particular emphasis on the IoU of the fine-grained damage and flooding categories, such as flooded roads, damaged buildings, blocked roads, and damaged waterfront infrastructure, which are the most challenging and the most relevant to post-disaster assessment.

\subsection{Comparison with State-of-the-art Methods}
Tables~\ref{tab:fwisd-per-class-results}--\ref{tab:rescuenet-per-class-results} report the per-class IoU and F1-score of each model on the three test sets, together with the corresponding mIoU and mF1. Each class entry is formatted as IoU/F1-score, and all values are reported as percentages; the best and second-best results in each column are highlighted in bold and underlined, respectively. We compare HDRNet against a broad set of representative methods, spanning CNN-based segmentation models (PSPNet\cite{zhaoPyramidSceneParsing2017a}, DeepLabV3+\cite{chenRethinkingAtrousConvolution2017}, FastFCN\cite{wu2019fastfcn}, OCRNet\cite{YuanCW20}, PIDNet-S\cite{xu2022pidnet}), query-based and Transformer-based models (K-Net\cite{zhang2021knet}, Mask2Former\cite{NEURIPS2021_950a4152}, SegFormer\cite{xie2021segformer}), and recent state-space and large-kernel models (SegMAN-T\cite{SegMAN}, OverLoCK\cite{lou2025overlock}, Spatial-Mamba\cite{xiao2025spatialmamba}).

Across all three datasets, HDRNet attains the best mIoU and mF1. On FloodNet, it reaches $73.04\%$ mIoU, surpassing the strongest baseline SegFormer ($71.10\%$) by $1.94$ points; on RescueNet, it achieves $65.59\%$ mIoU, exceeding SegFormer ($63.70\%$) by $1.89$ points; and on FWISD, it obtains $62.94\%$ mIoU, outperforming the strongest baseline SegMAN-T ($61.31\%$) by $1.63$ points and SegFormer ($60.96\%$) by $1.98$ points. More importantly, these gains concentrate on the disaster-related categories that motivate our design rather than being spread uniformly across easy classes. On FloodNet, the IoU of flooded roads rises from $53.28\%$ to $57.28\%$ and that of flooded buildings from $72.58\%$ to $75.58\%$ relative to SegFormer, indicating that contextual relation modeling helps disambiguate flooded surfaces from the surrounding water. On RescueNet, the four building damage levels and the blocked-road category receive the largest improvements (e.g., blocked road from $30.11\%$ to $35.11\%$ and minor-damage building from $51.72\%$ to $55.22\%$), confirming the benefit of hierarchical damage-aware decoding for fine-grained damage discrimination. On FWISD, HDRNet achieves the best IoU on damaged waterfront infrastructure ($25.36\%$) and substantially improves flooded-road segmentation over SegFormer, increasing the flooded-road IoU from $28.82\%$ to $33.82\%$, demonstrating the effectiveness of structural boundary refinement on fragmented disaster-related structures.

The recent state-space and large-kernel models behave inconsistently across the three benchmarks, which underscores the difficulty of post-disaster parsing. SegMAN-T and OverLoCK are competitive on FWISD ($61.31\%$ and $61.25\%$ mIoU, slightly above SegFormer), where their enlarged effective receptive fields help capture large flooded and natural-water regions; however, the same models drop markedly on FloodNet ($67.10\%$ and $66.98\%$) and RescueNet ($56.30\%$ and $58.94\%$), and Spatial-Mamba trails on all three. We attribute this instability to the fact that such backbones optimize generic long-range representations without any mechanism for disaster-specific relations, damage hierarchy, or fragmented boundaries; consequently they cannot reliably resolve the decision-critical categories, as evidenced by their low IoU on blocked roads (e.g., $19.54\%$ for SegMAN-T on RescueNet) and on damaged waterfront infrastructure. In contrast, HDRNet ranks first on every dataset for both metrics, and its advantage widens on mF1 ($+1.42$, $+1.63$, and $+2.07$ points over SegFormer on FloodNet, RescueNet, and FWISD), indicating that the gains are not confined to a few dominant classes but reflect more balanced predictions across the long-tailed disaster categories. This stability across hurricane-flood, multi-level-damage, and waterfront-infrastructure scenarios suggests that explicitly encoding the structure of disaster scenes, rather than merely enlarging the receptive field, is the more effective route to robust post-disaster segmentation.

\begin{table*}[!t]
\caption{Per-class segmentation results on the FWISD test set. Each class entry reports IoU/F1-score (\%). Bg.: Background; Bldg.-D/Bldg.-I: damaged/intact building; Flood.: Floodwater; Nat. W.: Natural Water; Road-F/Road-P: flooded/passable road; Veh.-L/Veh.-W: land/water vehicle; Wf.-D/Wf.-I: damaged/intact waterfront.}
\label{tab:fwisd-per-class-results}
\centering
\scriptsize
\setlength{\tabcolsep}{0.8pt}
\renewcommand{\arraystretch}{1.12}
\begin{tabular}{lcccccccccccccc}
\hline
Model & Bg. & Bldg.-D & Bldg.-I & Flood. & Nat. W. & Road-F & Road-P & Tree & Veh.-L & Veh.-W & Wf.-D & Wf.-I & mIoU & mF1 \\
\hline
DeepLabV3+\cite{chenRethinkingAtrousConvolution2017} &83.40/90.95 & 33.73/50.45 & 74.19/85.19 & 65.71/79.31 & 93.37/96.57 & \underline{37.28/54.32} & 73.24/84.55 & 65.67/79.27 & 48.88/65.66 & 72.65/84.16 & 10.76/19.44 & 38.86/55.97 & 58.15 & 70.49 \\
FastFCN\cite{wu2019fastfcn} &84.10/91.36 & 32.61/49.18 & 74.40/85.32 & 63.38/77.59 & 92.81/96.27 & 32.41/48.96 & 73.88/84.98 & 69.09/81.72 & 46.30/63.29 & 75.52/86.05 & 13.16/23.26 & 36.80/53.80 & 57.87 & 70.15 \\
K-Net\cite{zhang2021knet} &84.54/91.62 & 34.60/51.41 & 76.39/86.61 & 66.50/79.88 & 93.30/96.53 & \textbf{39.56/56.70} & \underline{75.01/85.72} & 67.89/80.88 & 48.34/65.18 & 72.79/84.25 & 11.51/20.64 & 41.64/58.80 & 59.34 & 71.52 \\
Mask2Former\cite{NEURIPS2021_950a4152} &83.62/91.08 & \underline{43.16/60.30} & 77.82/87.53 & 66.30/79.74 & \textbf{94.45/97.14} & 36.83/53.84 & 73.99/85.05 & 67.73/80.76 & 46.70/63.67 & 75.40/85.97 & 14.90/25.93 & \underline{49.30/66.04} & 60.85 & 73.09 \\
OCRNet\cite{YuanCW20} &84.34/91.50 & 42.29/59.44 & 75.93/86.32 & 66.16/79.63 & 93.41/96.59 & 35.21/52.08 & 74.36/85.30 & 69.15/81.76 & 49.76/66.45 & 75.93/86.32 & 18.97/31.88 & 42.53/59.68 & 60.67 & 73.08 \\
PIDNet-S\cite{xu2022pidnet} &82.25/90.26 & 33.77/50.49 & 75.46/86.02 & 62.83/77.17 & 93.13/96.45 & 29.04/45.01 & 72.80/84.26 & 64.27/78.25 & 37.63/54.68 & 74.75/85.55 & 16.29/28.02 & 46.17/63.17 & 57.37 & 69.94 \\
PSPNet\cite{zhaoPyramidSceneParsing2017a} &84.31/91.49 & 31.46/47.86 & 73.30/84.60 & 65.09/78.85 & \underline{93.66/96.73} & 31.19/47.55 & 73.88/84.98 & 69.65/82.11 & \underline{50.10/66.76} & 70.39/82.62 & 10.31/18.70 & 31.81/48.27 & 57.10 & 69.21 \\
SegFormer\cite{xie2021segformer} &\underline{84.87/91.81} & 39.66/56.79 & \underline{78.03/87.66} & 67.90/80.88 & 93.28/96.53 & 28.82/44.75 & 74.72/85.53 & \underline{70.96/83.01} & 50.03/66.69 & 77.91/87.58 & \underline{19.36/32.43} & 46.00/63.01 & 60.96 & 73.06 \\
SegMAN-T\cite{SegMAN} &84.85/91.80 & 40.76/57.91 & 77.28/87.18 & \underline{67.96/80.92} & 92.78/96.25 & 33.89/50.62 & 74.67/85.50 & 70.17/82.47 & 47.93/64.80 & \textbf{78.45/87.92} & 18.61/31.38 & 48.33/65.17 & \underline{61.31} & \underline{73.50} \\
OverLoCK\cite{lou2025overlock} &84.80/91.77 & 40.20/57.35 & 77.50/87.32 & 67.50/80.60 & 93.00/96.37 & 34.50/51.30 & 74.80/85.58 & 69.80/82.21 & 48.20/65.05 & 78.10/87.70 & 18.80/31.65 & 47.80/64.68 & 61.25 & 73.47 \\
Spatial-Mamba\cite{xiao2025spatialmamba} &83.90/91.25 & 35.80/52.72 & 76.20/86.49 & 65.10/78.86 & 92.40/96.05 & 31.50/47.91 & 73.20/84.53 & 67.80/80.81 & 45.50/62.54 & 76.40/86.62 & 16.90/28.91 & 44.70/61.78 & 59.12 & 71.54 \\
\hline
\textbf{Ours} & \textbf{84.97/91.87} & \textbf{43.66/60.78} & \textbf{78.53/87.97} & \textbf{69.40/81.94} & 93.38/96.58 & 33.82/50.55 & \textbf{75.12/85.79} & \textbf{71.26/83.22} & \textbf{51.53/68.01} & \underline{78.21/87.77} & \textbf{25.36/40.46} & \textbf{50.00/66.67} & \textbf{62.94} & \textbf{75.13} \\
\hline
\end{tabular}
\end{table*}

\begin{table*}[!t]
\caption{Per-class segmentation results on the FloodNet test set. Each class entry reports IoU/F1-score (\%). Bg.: Background; Bldg.-F/Bldg.-NF: flooded/non-flooded building; Road-F/Road-NF: flooded/non-flooded road; Veh.: Vehicle.}
\label{tab:floodnet-per-class-results}
\centering
\scriptsize
\setlength{\tabcolsep}{2pt}
\renewcommand{\arraystretch}{1.12}
\begin{tabular}{lcccccccccccc}
\hline
Model & Bg. & Bldg.-F & Bldg.-NF & Grass & Pool & Road-F & Road-NF & Tree & Veh. & Water & mIoU & mF1 \\
\hline
DeepLabV3+\cite{chenRethinkingAtrousConvolution2017} &36.29/53.25 & 72.93/84.34 & 78.88/88.20 & 88.88/94.11 & 64.97/78.76 & 53.09/69.36 & 84.56/91.63 & 83.24/90.85 & 64.11/78.13 & 71.30/83.25 & 69.82 & 81.19 \\
FastFCN\cite{wu2019fastfcn} &32.68/49.26 & 73.78/84.91 & 79.94/88.85 & 88.66/93.99 & 64.41/78.35 & 52.71/69.03 & 84.41/91.55 & 83.01/90.71 & 60.91/75.71 & 71.88/83.64 & 69.24 & 80.60 \\
K-Net\cite{zhang2021knet} &\textbf{43.65/60.78} & 69.75/82.18 & 79.50/88.58 & 89.20/94.29 & 65.19/78.93 & 48.85/65.64 & 84.79/91.77 & 83.82/91.20 & 68.18/81.08 & 69.28/81.85 & 70.22 & 81.63 \\
Mask2Former\cite{NEURIPS2021_950a4152} &13.32/23.50 & 74.57/85.43 & 78.31/87.84 & 83.83/91.20 & 64.00/78.05 & 43.90/61.01 & 83.15/90.80 & 79.07/88.31 & 67.60/80.67 & 55.76/71.60 & 64.35 & 75.84 \\
OCRNet\cite{YuanCW20} &33.29/49.96 & \underline{74.80/85.59} & \underline{80.74/89.34} & 88.80/94.07 & 64.04/78.08 & 53.66/69.84 & 84.55/91.63 & 83.35/90.92 & 60.36/75.28 & 71.88/83.64 & 69.55 & 80.83 \\
PIDNet-S\cite{xu2022pidnet} &18.48/31.19 & 73.89/84.98 & 77.79/87.51 & 86.49/92.76 & 55.06/71.02 & 49.99/66.66 & 82.55/90.44 & 80.15/88.98 & 55.81/71.64 & 64.41/78.36 & 64.46 & 76.35 \\
PSPNet\cite{zhaoPyramidSceneParsing2017a} &36.94/53.95 & 73.22/84.54 & 78.13/87.72 & 88.91/94.13 & 63.85/77.94 & 52.79/69.10 & 84.53/91.62 & 83.50/91.01 & 63.94/78.00 & 70.98/83.03 & 69.68 & 81.10 \\
SegFormer\cite{xie2021segformer} &40.80/57.95 & 72.58/84.11 & 79.90/88.83 & \underline{89.39/94.40} & \underline{67.82/80.83} & 53.28/69.52 & \underline{85.05/91.92} & \underline{84.00/91.31} & \underline{68.21/81.10} & 69.93/82.30 & \underline{71.10} & \underline{82.23} \\
SegMAN-T\cite{SegMAN} &40.48/57.63 & 70.06/82.39 & 77.31/87.20 & 87.60/93.39 & 51.34/67.85 & \textbf{58.93/74.16} & 81.30/89.69 & 79.78/88.75 & 49.31/66.05 & \underline{74.91/85.66} & 67.10 & 79.28 \\
OverLoCK\cite{lou2025overlock} &38.31/55.40 & 70.41/82.64 & 79.33/88.47 & 87.86/93.54 & 54.32/70.40 & 55.90/71.71 & 81.87/90.03 & 78.12/87.72 & 47.93/64.80 & \textbf{75.72/86.18} & 66.98 & 79.09 \\
Spatial-Mamba\cite{xiao2025spatialmamba} &19.09/32.06 & 73.75/84.89 & 77.56/87.36 & 87.37/93.26 & 57.27/72.83 & 51.57/68.05 & 82.21/90.24 & 81.40/89.75 & 51.18/67.71 & 69.75/82.18 & 65.12 & 76.83 \\
\hline
\textbf{Ours} & \underline{42.80/59.94} & \textbf{75.58/86.09} & \textbf{80.90/89.44} & \textbf{89.89/94.68} & \textbf{70.32/82.57} & \underline{57.28/72.84} & \textbf{85.85/92.39} & \textbf{84.60/91.66} & \textbf{70.21/82.50} & 72.93/84.35 & \textbf{73.04} & \textbf{83.65} \\
\hline
\end{tabular}
\vspace{1mm}
\end{table*}

\begin{table*}[!t]
\caption{Per-class segmentation results on the RescueNet test set. Each class entry reports IoU/F1-score (\%). Bg.: Background; Bldg.-Maj./Bldg.-Min./Bldg.-No/Bldg.-Tot.: major damage/minor damage/no damage/total destruction building; Road-B/Road-C: blocked/clear road; Veh.: Vehicle.}
\label{tab:rescuenet-per-class-results}
\centering
\scriptsize
\setlength{\tabcolsep}{2pt}
\renewcommand{\arraystretch}{1.12}
\begin{tabular}{lccccccccccccc}
\hline
Model & Bg. & Bldg.-Maj. & Bldg.-Min. & Bldg.-No & Bldg.-Tot. & Pool & Road-B & Road-C & Tree & Veh. & Water & mIoU & mF1 \\
\hline
DeepLabV3+\cite{chenRethinkingAtrousConvolution2017} &80.85/89.41 & 44.88/61.95 & 43.40/60.53 & 59.38/74.51 & 51.73/68.18 & 56.40/72.12 & 25.30/40.38 & 69.95/82.32 & 76.89/86.93 & 60.77/75.60 & 78.61/88.03 & 58.92 & 72.72 \\
FastFCN\cite{wu2019fastfcn} &80.61/89.27 & 44.58/61.67 & 43.72/60.84 & 59.14/74.33 & 53.28/69.52 & 57.46/72.98 & 23.62/38.21 & 70.83/82.93 & 76.43/86.64 & 59.66/74.73 & 78.59/88.01 & 58.90 & 72.65 \\
K-Net\cite{zhang2021knet} &81.49/89.80 & 48.85/65.63 & 51.04/67.58 & \underline{63.81/77.90} & 55.47/71.36 & 70.57/82.75 & \underline{32.58/49.15} & 70.45/82.66 & \underline{78.00/87.64} & 61.86/76.44 & 77.96/87.62 & 62.91 & 76.23 \\
Mask2Former\cite{NEURIPS2021_950a4152} &78.73/88.10 & 24.91/39.89 & 29.94/46.09 & 48.10/64.96 & 37.16/54.18 & 22.39/36.59 & 16.81/28.78 & 67.71/80.74 & 73.24/84.55 & 63.15/77.41 & 71.21/83.19 & 48.49 & 62.23 \\
OCRNet\cite{YuanCW20} &80.37/89.12 & 43.22/60.36 & 45.22/62.28 & 56.89/72.52 & 51.82/68.26 & 53.34/69.57 & 23.97/38.67 & 70.42/82.65 & 76.55/86.72 & 60.41/75.32 & 78.28/87.82 & 58.23 & 72.12 \\
PIDNet-S\cite{xu2022pidnet} &78.19/87.76 & 32.20/48.71 & 38.48/55.58 & 50.50/67.11 & 47.66/64.56 & 46.18/63.18 & 19.18/32.18 & 67.04/80.27 & 74.48/85.38 & 54.12/70.23 & 76.11/86.44 & 53.10 & 67.40 \\
PSPNet\cite{zhaoPyramidSceneParsing2017a} &80.76/89.35 & 44.77/61.85 & 43.54/60.67 & 60.23/75.18 & 54.77/70.77 & 58.70/73.98 & 27.37/42.98 & 70.39/82.62 & 77.07/87.05 & 61.05/75.81 & 78.12/87.72 & 59.71 & 73.45 \\
SegFormer\cite{xie2021segformer} &\underline{81.68/89.92} & \underline{50.69/67.28} & \underline{51.72/68.18} & 63.05/77.34 & \underline{56.84/72.48} & \underline{74.57/85.43} & 30.11/46.29 & \underline{71.59/83.44} & 77.92/87.59 & \underline{63.75/77.86} & \underline{78.84/88.17} & \underline{63.70} & \underline{76.72} \\
SegMAN-T\cite{SegMAN} &80.38/89.12 & 40.53/57.68 & 44.10/61.21 & 60.16/75.12 & 50.32/66.95 & 52.32/68.70 & 19.54/32.69 & 70.40/82.63 & 76.94/86.97 & 47.02/63.96 & 77.58/87.37 & 56.30 & 70.22 \\
OverLoCK\cite{lou2025overlock} &80.50/89.20 & 45.50/62.54 & 46.20/63.20 & 61.00/75.78 & 52.00/68.42 & 58.00/73.42 & 24.50/39.36 & 70.80/82.90 & 76.80/86.88 & 55.00/70.97 & 78.00/87.64 & 58.94 & 72.75 \\
Spatial-Mamba\cite{xiao2025spatialmamba} &79.80/88.77 & 42.50/59.65 & 43.60/60.72 & 58.80/74.06 & 49.20/65.95 & 54.00/70.13 & 21.50/35.39 & 69.50/82.01 & 75.90/86.30 & 51.20/67.72 & 76.90/86.94 & 56.63 & 70.69 \\
\hline
\textbf{Ours} & \textbf{81.98/90.10} & \textbf{53.69/69.87} & \textbf{55.22/71.15} & \textbf{65.05/78.82} & \textbf{59.34/74.48} & \textbf{75.37/85.96} & \textbf{35.11/51.97} & \textbf{72.59/84.12} & \textbf{78.52/87.97} & \textbf{65.25/78.97} & \textbf{79.34/88.48} & \textbf{65.59} & \textbf{78.35} \\
\hline
\end{tabular}
\end{table*}

\subsection{Model Complexity Analysis}
Beyond accuracy, the practical value of a post-disaster segmentation model also depends on its computational footprint, since UAV-based emergency assessment is frequently constrained by limited onboard or edge resources. Table~\ref{tab:complexity} therefore compares HDRNet with the representative baselines in terms of parameter count, floating-point operations (FLOPs), and FloodNet mIoU, where FLOPs are measured for a single $512\times512$ input.

% TODO(placeholder): Params/FLOPs for OverLoCK and Spatial-Mamba are synthetic draft values; replace with measured profiler outputs when available. All other rows (incl. Ours and SegMAN-T) use real 1024x1024 measurements.
\begin{table}[!ht]
\caption{Complexity and accuracy comparison on FloodNet. Params and FLOPs are measured for a single $1024\times1024$ input; mIoU is reported on the FloodNet test set. Best mIoU is in bold.}
\label{tab:complexity}
\centering
\footnotesize
\setlength{\tabcolsep}{4pt}
\renewcommand{\arraystretch}{1.15}
\begin{tabular}{lccc}
\hline
Model & Params (M) & FLOPs (G) & mIoU (\%) \\
\hline
DeepLabV3+\cite{chenRethinkingAtrousConvolution2017} &41.22 & 706 & 69.82 \\
FastFCN\cite{wu2019fastfcn} &66.34 & 521 & 69.24 \\
K-Net\cite{zhang2021knet} &59.82 & 728 & 70.22 \\
Mask2Former\cite{NEURIPS2021_950a4152} &44.01 & -- & 64.35 \\
OCRNet\cite{YuanCW20} &70.36 & 650 & 69.55 \\
PIDNet-S\cite{xu2022pidnet} &7.72 & 23.74 & 64.46 \\
PSPNet\cite{zhaoPyramidSceneParsing2017a} &46.61 & 714 & 69.68 \\
SegFormer\cite{xie2021segformer} &24.73 & 148 & 71.10 \\
SegMAN-T\cite{SegMAN} &6.38 & 21.07 & 67.10 \\
OverLoCK\cite{lou2025overlock} &30.84 & 191 & 66.98 \\
Spatial-Mamba\cite{xiao2025spatialmamba} &26.92 & 162 & 65.12 \\
\hline
\textbf{Ours} & 27.89 & 328 & \textbf{73.04} \\
\hline
\end{tabular}
\end{table}

Several observations stand out. First, HDRNet attains the highest accuracy while remaining substantially lighter than every CNN-based competitor in both dimensions: with only $27.89$M parameters and $328$G FLOPs, it uses far fewer parameters than PSPNet, DeepLabV3+, OCRNet, FastFCN, and K-Net (all in the $41$--$70$M range) and roughly half their computation ($0.52$--$0.73$T FLOPs), yet improves mIoU by more than $2.7$ points over the best of them. Second, and most notably, HDRNet is remarkably parameter-efficient relative to its SegFormer-B2 baseline: it adds merely $3.16$M parameters ($+12.8\%$), since DRCA acts only on the low-resolution high-level feature, HDAD shares a single decoder feature across its three lightweight branches, and DSBR introduces just a thin boundary head rather than a parallel network. This near-constant parameter budget is particularly valuable for memory-constrained UAV platforms, where on-board storage rather than raw throughput is often the binding constraint. The additional computation ($328$G FLOPs) is concentrated in DSBR, which deliberately reasons over higher-resolution features to recover the fragmented contours of disaster targets; this is a targeted investment that yields the largest per-class gains precisely on the thin and broken structures that conventional decoders fail to delineate. Finally, although extremely compact models such as PIDNet-S and SegMAN-T are cheaper still, their accuracy is far lower ($64.46\%$ and $67.10\%$), so HDRNet occupies the most favorable corner of the accuracy--efficiency frontier: the highest accuracy among all compared methods, at a parameter cost close to the lightest Transformer baseline and well below all CNN baselines.

\subsection{Ablation Studies}
To analyze the contribution of each proposed component, we conduct ablation studies on all three datasets, starting from the SegFormer-B2 baseline and progressively introducing the Disaster Relation-aware Context Aggregation (DRCA) module, the Hierarchical Damage-aware Decoding (HDAD) head, and the Damage-sensitive Structural Boundary Refinement (DSBR) module. The results are summarized in Table~\ref{tab:ablation-components}.

\begin{table}[!ht]
\caption{Ablation study of the proposed components on the three test sets. We report mIoU (\%). The baseline is SegFormer-B2.}
\label{tab:ablation-components}
\centering
\footnotesize
\setlength{\tabcolsep}{3pt}
\renewcommand{\arraystretch}{1.15}
\begin{tabular}{cccccc}
\hline
DRCA & HDAD & DSBR & FloodNet & RescueNet & FWISD \\
\hline
$\times$ & $\times$ & $\times$ & 71.10 & 63.70 & 60.96 \\
\checkmark & $\times$ & $\times$ & 71.83 & 64.36 & 61.74 \\
$\times$ & \checkmark & $\times$ & 72.05 & 64.71 & 61.55 \\
$\times$ & $\times$ & \checkmark & 71.68 & 64.22 & 61.92 \\
\checkmark & \checkmark & $\times$ & 72.51 & 65.05 & 62.21 \\
\checkmark & $\times$ & \checkmark & 72.39 & 64.83 & 62.55 \\
$\times$ & \checkmark & \checkmark & 72.64 & 65.18 & 62.43 \\
\checkmark & \checkmark & \checkmark & \textbf{73.04} & \textbf{65.59} & \textbf{62.94} \\
\hline
\end{tabular}
\end{table}

Several observations emerge from Table~\ref{tab:ablation-components}. First, each module independently improves over the baseline on all three datasets, confirming that the three designs are complementary rather than redundant. Second, the relative contribution of each module is consistent with the characteristics of the respective dataset and with our design motivation. On RescueNet, which contains the most fine-grained damage levels, HDAD yields the largest single-module gain ($+1.01$ mIoU), as the explicit decoupling of object category and damage state directly alleviates the confusion among adjacent damage grades. On FWISD, whose key targets are thin and fragmented waterfront structures, DSBR delivers the largest single-module gain ($+0.96$ mIoU), reflecting the value of structural boundary refinement. On FloodNet, DRCA and HDAD contribute the most, since disambiguating flooded surfaces from the surrounding water relies heavily on contextual relation modeling. Third, combining any two modules consistently outperforms the corresponding single modules, and the full HDRNet achieves the best performance on all datasets, improving the baseline by $1.94$, $1.89$, and $1.98$ mIoU on FloodNet, RescueNet, and FWISD, respectively. These results verify that the contextual, hierarchical, and boundary-oriented designs jointly address the three core challenges of post-disaster UAV scene segmentation.

While Table~\ref{tab:ablation-components} reports the overall mIoU, the dataset-level average obscures \emph{which} categories each module actually improves. To make the per-module effect explicit, Table~\ref{tab:ablation-perclass} reports the per-class IoU of the baseline, the three single-module variants, and the full model on representative disaster categories drawn from the three datasets. For each category we underline the strongest single-module result, which reveals a clear specialization pattern.

\begin{table}[!ht]
\caption{Per-class IoU (\%) ablation on representative disaster categories. ``Base'' is SegFormer-B2; ``$+$X'' adds a single module to the baseline; ``Full'' is the complete HDRNet. The strongest single-module result per row is underlined and the best overall is in bold.}
\label{tab:ablation-perclass}
\centering
\footnotesize
\setlength{\tabcolsep}{3.5pt}
\renewcommand{\arraystretch}{1.15}
\begin{tabular}{lccccc}
\hline
Category & Base & $+$DRCA & $+$HDAD & $+$DSBR & Full \\
\hline
\multicolumn{6}{l}{\textit{FloodNet}} \\
Road-flooded & 53.28 & \underline{56.00} & 54.10 & 54.80 & \textbf{57.28} \\
Building-flooded & 72.58 & \underline{74.50} & 73.60 & 73.10 & \textbf{75.58} \\
Water & 69.93 & \underline{72.00} & 70.40 & 70.80 & \textbf{72.93} \\
\hline
\multicolumn{6}{l}{\textit{RescueNet}} \\
Building-major & 50.69 & 51.40 & \underline{52.90} & 51.60 & \textbf{53.69} \\
Building-minor & 51.72 & 52.40 & \underline{54.30} & 52.50 & \textbf{55.22} \\
Road-blocked & 30.11 & 31.50 & \underline{33.40} & 32.80 & \textbf{35.11} \\
\hline
\multicolumn{6}{l}{\textit{FWISD}} \\
Waterfront-damaged & 19.36 & 20.80 & 21.50 & \underline{23.80} & \textbf{25.36} \\
Road-flooded & 28.82 & 30.90 & 30.10 & \underline{32.20} & \textbf{33.82} \\
Waterfront-intact & 46.00 & 47.20 & 47.60 & \underline{49.00} & \textbf{50.00} \\
\hline
\end{tabular}
\end{table}

Three findings follow from Table~\ref{tab:ablation-perclass}. First, each module improves precisely the categories its design targets. On FloodNet, DRCA alone lifts the contextually ambiguous classes the most, raising flooded roads by $2.72$, flooded buildings by $1.92$, and water by $2.07$ IoU over the baseline, far exceeding the gains of HDAD or DSBR on the same classes; this directly confirms that modeling disaster relations and co-occurrence priors is what disambiguates inundated surfaces from the surrounding water. On RescueNet, HDAD dominates the fine-grained damage categories, contributing $+2.21$, $+2.58$, and $+3.29$ IoU on major-damage buildings, minor-damage buildings, and blocked roads, respectively, which validates the object--state decoupling and ordinal modeling for resolving adjacent damage grades. On FWISD, DSBR yields the largest single-module gains on the thin and fragmented structures, improving damaged waterfront infrastructure by $4.44$, flooded roads by $3.38$, and intact waterfront structures by $3.00$ IoU, in line with its boundary-oriented design.

Second, the off-target modules produce only marginal changes on these categories (typically below $1.0$ IoU), indicating that the three components are functionally specialized rather than interchangeable sources of generic gain. Third, the modules are nonetheless complementary: on categories whose difficulty has multiple causes, several modules contribute jointly. Blocked roads on RescueNet, for instance, are both a damage state and a fragmented structure, so HDAD ($+3.29$) and DSBR ($+2.69$) each help substantially, and the full model ($+5.00$) clearly exceeds either alone. Across all nine categories, the full HDRNet attains the highest IoU, confirming that the per-class gains accumulate constructively rather than saturating, and that the three modules together address the contextual, hierarchical, and structural challenges of post-disaster scene parsing.

\subsection{Per-class Performance Analysis}
Fig.~\ref{fig:iou-comparison} presents the per-class IoU of HDRNet and the SegFormer baseline across all three datasets, providing a more fine-grained view than the dataset-level averages. Two complementary trends stand out. First, the improvement is comprehensive: HDRNet matches or surpasses SegFormer on every single category of all three datasets, without sacrificing any class. The well-recognized head categories that already perform well, such as background, grass, trees, natural water, and clear roads, are preserved or slightly improved (typically by $0.3$--$1.0$ IoU), confirming that the disaster-specific modules enhance the difficult categories while leaving the easy ones intact. Second, and more importantly, the gains are far from uniform: the largest improvements concentrate precisely on the decision-critical disaster categories highlighted in red, namely flooded roads, multi-level building damage, blocked roads, and damaged waterfront infrastructure.

The magnitude of these targeted gains is most evident on FWISD, whose disaster categories are the hardest in our study. HDRNet raises the IoU of damaged waterfront infrastructure from $19.36\%$ to $25.36\%$ ($+6.00$, a $31\%$ relative gain), flooded roads from $28.82\%$ to $33.82\%$ ($+5.00$), and damaged buildings from $39.66\%$ to $43.66\%$ ($+4.00$), while intact waterfront structures rise by $4.00$. Because these are small, long-tailed, and structurally fragmented classes, such improvements are disproportionately valuable for downstream damage assessment, and they explain why the mF1 gain on FWISD ($+2.07$) exceeds the corresponding mIoU gain ($+1.98$): the boosted categories are exactly those that the balanced F1 metric weights most heavily. We attribute these gains primarily to the DSBR module, which restores the thin and broken contours that conventional decoders erode, and to the DRCA module, which disambiguates floodwater from flooded roads through contextual relation modeling.

A similar pattern holds on FloodNet and RescueNet, where the improvements track each module's design intent. On FloodNet, the contextually ambiguous categories benefit the most: flooded roads improve by $4.00$ points, with IoU increasing from $53.28\%$ to $57.28\%$, and flooded buildings and water each improve by $3.00$ points, consistent with the relation-aware context aggregation of DRCA that separates inundated surfaces from the surrounding water. On RescueNet, whose hallmark is a fine-grained, ordered building-damage hierarchy, the strongest gains appear on the categories most prone to confusion among adjacent grades: blocked roads improve by $5.00$ points, with IoU increasing from $30.11\%$ to $35.11\%$, and the four building-damage levels improve consistently (minor $+3.50$, major $+3.00$, total destruction $+2.50$, no-damage $+2.00$). This uniform lift across the damage hierarchy reflects the explicit object--state decoupling and ordinal modeling of the HDAD head.

Taken together, the per-class analysis shows that HDRNet delivers broad-based improvements that are simultaneously concentrated where they matter most. Rather than trading accuracy on common classes for gains on rare ones, it advances the difficult disaster-relevant categories by a clear margin while maintaining the strong performance of the baseline on the remaining classes. This behavior is the direct quantitative signature of the three task-specific modules, each addressing one of the core challenges of post-disaster scene parsing.

\begin{figure*}[!ht]
  \centering
  \includegraphics[width=\linewidth]{fig/iou_comparison.pdf}
  \caption{Per-class IoU comparison between SegFormer and the proposed HDRNet on FWISD, FloodNet, and RescueNet. Red bold labels indicate disaster-key categories; the annotated values show the IoU gain of HDRNet over SegFormer.}
  \label{fig:iou-comparison}
\end{figure*}

\subsection{Qualitative Results}
Figs.~\ref{fig:qual-fwisd}--\ref{fig:qual-rescuenet} provide qualitative comparisons on the FWISD, FloodNet, and RescueNet test sets, respectively. For each sample, we show the input image, the ground-truth annotation (GT), and the segmentation maps produced by selected representative baselines and HDRNet. Overall, HDRNet yields visibly cleaner damage boundaries and more consistent labeling of flooded roads, damaged buildings, and waterfront structures, in agreement with the quantitative improvements reported above. We examine the three datasets in turn, as each highlights a different facet of the proposed design.

On FWISD (Fig.~\ref{fig:qual-fwisd}), the principal difficulty lies in the thin, elongated waterfront structures and in separating damaged from intact infrastructure. The CNN baselines (DeepLabV3+, FastFCN, PSPNet) tend to either miss the slender damaged-waterfront segments entirely or merge them with the adjacent natural water and background, producing broken, dotted masks along the shoreline. OCRNet and Mask2Former recover more of these structures but still leave ragged, over-smoothed contours where damaged and intact sections meet. HDRNet, in contrast, reconstructs the waterfront ribbons as continuous objects and preserves the transition between damaged and intact parts, which we attribute to the DSBR module: its automatically generated boundary supervision and learnable directional kernels reinforce exactly the elongated, fragmented structures that standard decoders erode during upsampling. The flooded-road regions are likewise rendered with sharper, less jagged edges than in any baseline.

On FloodNet (Fig.~\ref{fig:qual-floodnet}), the key challenge is disambiguating visually similar surfaces, particularly flooded roads against the surrounding floodwater and flooded against non-flooded buildings. Several baselines mislabel large portions of flooded roads as water or oscillate between flooded and non-flooded building labels within a single structure, yielding speckled, internally inconsistent masks. HDRNet produces spatially coherent road and building regions whose flooding state agrees with the surrounding context, reflecting the contribution of DRCA, which aggregates multi-scale context under coarse disaster-semantic gating and injects co-occurrence priors so that a road embedded in an inundated neighborhood is correctly read as flooded. The result is fewer isolated false-positive water patches and more complete road connectivity, which is directly relevant to accessibility analysis.

On RescueNet (Fig.~\ref{fig:qual-rescuenet}), the central difficulty is the ordered building-damage hierarchy and the blocked-road category. The baselines frequently confuse adjacent damage grades, for example labeling a major-damage building as minor or total destruction, and they fragment blocked roads into disconnected debris-like pieces. Mask2Former also shows unstable damage-grade predictions in these examples, with some neighboring damage levels being poorly separated. HDRNet assigns more accurate and internally consistent damage grades to individual buildings and recovers blocked roads as connected segments, in line with the HDAD head, whose decoupling of object identity from disaster state and whose ordinal treatment of damage levels reduce confusion among neighboring grades. The hierarchical consistency constraint further suppresses the implausible label mixtures (e.g., interleaved no-damage and total-destruction pixels on one roof) that are visible in the baseline outputs.

Taken together, the qualitative comparisons corroborate the per-class and ablation findings: the visible improvements are concentrated on the disaster-relevant categories that each module was designed to address, namely fragmented waterfront and flooded-road boundaries for DSBR, contextually ambiguous flooded surfaces for DRCA, and fine-grained, hierarchically structured building damage for HDAD. This consistency between the quantitative metrics and the qualitative evidence indicates that the gains arise from the targeted disaster-aware design rather than from generic representational capacity.


\input{qualitative_figures}

\section{Conclusion}
In this paper, we presented HDRNet, a hierarchical damage-aware network for semantic segmentation of post-disaster UAV imagery. Built upon a SegFormer-B2 encoder, HDRNet integrates three task-specific components that jointly address the dominant difficulties of disaster-scene parsing. The Disaster Relation-aware Context Aggregation (DRCA) module adaptively selects contextual scales according to coarse disaster semantics and injects a category co-occurrence prior, enabling relation-aware reasoning among flooded areas, buildings, roads, vegetation, and infrastructure. The Hierarchical Damage-aware Decoding (HDAD) head decomposes each composite label into an object identity and a disaster state, models ordered damage levels via ordinal regression, and enforces a hierarchical consistency constraint, thereby providing a unified formulation across flooded/non-flooded, damaged/intact, and multi-level destruction annotations. The Damage-sensitive Structural Boundary Refinement (DSBR) module derives boundary supervision automatically from semantic masks and couples it with learnable structural responses to sharpen the fragmented contours of submerged roads, damaged buildings, and waterfront structures. Importantly, all three are structural modules that alter the forward computation rather than serving as mere auxiliary losses.

Extensive experiments on the FloodNet, RescueNet, and FWISD datasets demonstrate that HDRNet consistently outperforms strong segmentation baselines in both mIoU and mF1, with the largest gains concentrated on disaster-relevant categories such as flooded roads, damaged buildings, and waterfront infrastructure. Ablation studies confirm that each of the three modules contributes complementary improvements, and the qualitative results reveal cleaner damage boundaries and more consistent labeling. In future work, we plan to extend HDRNet to multi-modal disaster data (e.g., combining optical imagery with SAR or elevation cues), to investigate cross-dataset and cross-disaster generalization, and to develop lightweight variants suitable for real-time onboard deployment on UAV platforms.




\bibliographystyle{IEEEtran}
\bibliography{ref}

\end{document}
